
\documentclass[11pt,letterpaper,oneside]{report}
\usepackage[T1]{fontenc}
\usepackage[utf8]{inputenc}
\usepackage{lmodern}
\usepackage{inconsolata}
\usepackage{geometry}
\geometry{margin=0.89in,headheight=16pt,headsep=19pt,footskip=26pt}
\usepackage{xcolor}
\definecolor{navy}{HTML}{173758}
\definecolor{teal}{HTML}{157B89}
\definecolor{pale}{HTML}{E9F0F3}
\definecolor{codegray}{HTML}{F5F7F9}
\definecolor{muted}{HTML}{667785}
\usepackage{booktabs,longtable,tabularx,array,calc}
\usepackage{amsmath,amssymb}
\usepackage{graphicx}
\usepackage{microtype}
\usepackage{enumitem}
\usepackage{fancyhdr}
\usepackage{titlesec}
\usepackage{listings}
\usepackage{upquote}
\usepackage{xurl}
\usepackage{hyperref}
\hypersetup{colorlinks=true,linkcolor=navy,urlcolor=teal,bookmarksopen=true,pdfauthor={Software Build Engineering Documentation Suite}}
\urlstyle{same}
\setcounter{tocdepth}{2}
\setcounter{secnumdepth}{2}
\setlength{\parindent}{0pt}
\setlength{\parskip}{5pt plus 1pt}
\setlength{\emergencystretch}{3em}
\setlist[itemize]{leftmargin=*,topsep=4pt,itemsep=2pt}
\setlist[enumerate]{leftmargin=*,topsep=4pt,itemsep=2pt}
\titleformat{\chapter}[display]{\normalfont\bfseries\color{navy}}{\fontsize{11}{12}\selectfont CHAPTER \thechapter}{7pt}{\fontsize{23}{28}\selectfont}
\titlespacing*{\chapter}{0pt}{-6pt}{19pt}
\titleformat{\section}{\Large\bfseries\color{navy}}{\thesection}{0.65em}{}
\titleformat{\subsection}{\large\bfseries\color{teal}}{\thesubsection}{0.65em}{}
\pagestyle{fancy}
\fancyhf{}
\fancyhead[L]{\small\textcolor{muted}{PROFESSIONAL CMAKE ENGINEERING}}
\fancyhead[R]{\small\textcolor{muted}{\nouppercase{\leftmark}}}
\fancyfoot[C]{\small\textcolor{muted}{\thepage}}
\renewcommand{\headrulewidth}{0.3pt}
\fancypagestyle{plain}{\fancyhf{}\fancyfoot[C]{\small\thepage}\renewcommand{\headrulewidth}{0pt}}
\lstdefinelanguage{cmake}{sensitive=false,morekeywords={cmake_minimum_required,project,add_library,add_executable,target_compile_features,target_include_directories,target_link_libraries,target_sources,target_compile_options,target_compile_definitions,target_link_options,add_subdirectory,include,if,else,elseif,endif,foreach,endforeach,function,endfunction,macro,endmacro,set,option,find_package,find_program,add_custom_command,add_custom_target,install,configure_file,configure_package_config_file,write_basic_package_version_file,set_target_properties,set_tests_properties,add_test,message,FetchContent_Declare,FetchContent_MakeAvailable,return,cmake_parse_arguments},morecomment=[l]{\#},morestring=[b]"}
\lstdefinelanguage{sh}{morekeywords={cmake,ctest,cpack,find,sha256sum,clang-tidy,echo,cp,make,ninja,git},morecomment=[l]{\#},morestring=[b]"}
\lstdefinelanguage{yaml}{morecomment=[l]{\#},morestring=[b]",morekeywords={name,on,permissions,jobs,steps,uses,run,with,env,matrix,include,if}}
\lstdefinelanguage{text}{}
\lstset{basicstyle=\ttfamily\footnotesize,keywordstyle=\color{navy}\bfseries,commentstyle=\color{muted},stringstyle=\color{teal},backgroundcolor=\color{codegray},frame=single,rulecolor=\color{pale},breaklines=true,breakatwhitespace=false,showstringspaces=false,columns=fullflexible,keepspaces=true,tabsize=2,xleftmargin=5pt,xrightmargin=5pt,aboveskip=9pt,belowskip=9pt}
\providecommand{\tightlist}{\setlength{\itemsep}{0pt}\setlength{\parskip}{0pt}}
\providecommand{\passthrough}[1]{#1}
\providecommand{\real}[1]{#1}

\begin{document}
\hypersetup{pageanchor=false}\begin{titlepage}
\color{navy}
\vspace*{1.15in}
{\Large\bfseries PROFESSIONAL CMAKE\\[7pt]SOFTWARE BUILD ENGINEERING\par}
\vspace{0.32in}\noindent\textcolor{teal}{\rule{0.9\linewidth}{2.3pt}}\par
\vspace{0.45in}
{\parbox{0.95\textwidth}{\raggedright\fontsize{24}{29}\selectfont\bfseries Modern CMake Implementation Handbook\par}}
\vspace{0.3in}
{\Large\color{teal}Runnable multi-target reference project\par}
\vfill
{\large DOCUMENT 03 OF 08\par}
\vspace{0.1in}
{\small Companion synthesis to Craig Scott, \emph{Professional CMake: A Practical Guide}, 20th Edition (2025).\\
Original learning/engineering material. Not affiliated with or endorsed by the author.\par}
\vspace{0.2in}
{\small 17 October 2026\par}
\end{titlepage}
\hypersetup{pageanchor=true}\clearpage\pagenumbering{roman}\tableofcontents\clearpage\pagenumbering{arabic}

\textbf{Goal.} Deliver a multi-target C++ project that builds, tests,
installs, exports a package, and can be consumed from an independent
source directory. \textbf{Working artifact:}
\passthrough{\lstinline!examples/reference-project!} in this suite.
\textbf{Provenance:} {[}BOOK-GROUNDED{]} target construction, CTest,
installation, exports and presets relate to Chapters 2--5, 16, 19--20,
27, 34--35, 40 and 42. {[}ENGINEERING EXTENSION{]} quality gates,
reproducibility checks and the reference implementation are new.

\chapter{Initial verification}\label{initial-verification}

\begin{lstlisting}[language=sh]
cmake --version
cmake -G  # list available generators
c++ --version
\end{lstlisting}

Use CMake 3.25 or newer, a functional C++17 compiler and Ninja where
possible. The demo uses standard-library facilities only; it does not
download dependencies. Check that a compiler change triggers a
\textbf{new build directory}, not a silent reuse of an existing cache.

\chapter{Repository layout and component
ownership}\label{repository-layout-and-component-ownership}

\begin{lstlisting}
reference-project/
  CMakeLists.txt
  CMakePresets.json
  cmake/BuildSuiteConfig.cmake.in
  include/bsuite/math.hpp
  src/math.cpp
  src/main.cpp
  tests/math_tests.cpp
  consumer/CMakeLists.txt
  consumer/main.cpp
  .github/workflows/build.yml
\end{lstlisting}

\begin{longtable}[]{@{}
  >{\raggedright\arraybackslash}p{(\columnwidth - 4\tabcolsep) * \real{0.3333}}
  >{\raggedright\arraybackslash}p{(\columnwidth - 4\tabcolsep) * \real{0.3333}}
  >{\raggedright\arraybackslash}p{(\columnwidth - 4\tabcolsep) * \real{0.3333}}@{}}
\toprule\noalign{}
\begin{minipage}[b]{\linewidth}\raggedright
Component
\end{minipage} & \begin{minipage}[b]{\linewidth}\raggedright
Responsibility
\end{minipage} & \begin{minipage}[b]{\linewidth}\raggedright
Consumer contract
\end{minipage} \\
\midrule\noalign{}
\endhead
\bottomrule\noalign{}
\endlastfoot
\passthrough{\lstinline!bsuite\_core!} & public C++ library, clean
header boundary & \passthrough{\lstinline!BuildSuite::core!} \\
\passthrough{\lstinline!bsuite\_cli!} & private CLI behavior and output
& executable artifact \\
\passthrough{\lstinline!bsuite\_tests!} & direct assertions on library
behavior & CTest test \passthrough{\lstinline!core.math!} \\
\passthrough{\lstinline!BuildSuiteConfig.cmake!} & imports installed
target definitions &
\passthrough{\lstinline!find\_package(BuildSuite CONFIG REQUIRED)!} \\
\passthrough{\lstinline!consumer/!} & independent install validation &
link only to \passthrough{\lstinline!BuildSuite::core!} \\
\end{longtable}

\chapter{Stage 1: Configure and build a
target}\label{stage-1-configure-and-build-a-target}

\begin{lstlisting}[language=sh]
cd examples/reference-project
cmake --preset ci-debug
cmake --build --preset ci-debug -j 4
\end{lstlisting}

The top-level \passthrough{\lstinline!CMakeLists.txt!} declares a named
target with a public \passthrough{\lstinline!FILE\_SET HEADERS!}, an
alias, a CLI and test executable. Public compile features are exported
as part of the library's consumer contract. The header file set avoids
hand-maintaining a second installation list. Look at the actual source
files for authoritative code: the code fragments in this handbook are
conceptual illustrations.

\begin{lstlisting}
add_library(bsuite_core STATIC src/math.cpp)
add_library(BuildSuite::core ALIAS bsuite_core)
target_compile_features(bsuite_core PUBLIC cxx_std_17)
target_sources(bsuite_core PUBLIC FILE_SET HEADERS
  BASE_DIRS "${CMAKE_CURRENT_SOURCE_DIR}/include"
  FILES "${CMAKE_CURRENT_SOURCE_DIR}/include/bsuite/math.hpp")
\end{lstlisting}

\textbf{Verification:}
\passthrough{\lstinline!cmake --build --preset ci-debug --verbose!}
should reveal a single compilation of
\passthrough{\lstinline!math.cpp!}; public includes must be available to
the CLI and unit tests via the link graph. Do not supply the include
path separately to each consumer.

\chapter{Stage 2: Define private and propagated
requirements}\label{stage-2-define-private-and-propagated-requirements}

\begin{lstlisting}
target_link_libraries(bsuite_cli PRIVATE BuildSuite::core)
target_link_libraries(bsuite_tests PRIVATE BuildSuite::core)
\end{lstlisting}

The CLI and test source compilation inherit the library's public include
directories and required C++ standard. These properties do \textbf{not}
imply that a runtime shared library is available on all systems. For a
shared library, also model symbol exports, relocatable loaders and ABI
compatibility.

\section{When to use each target
type}\label{when-to-use-each-target-type}

\begin{longtable}[]{@{}
  >{\raggedright\arraybackslash}p{(\columnwidth - 4\tabcolsep) * \real{0.3333}}
  >{\raggedright\arraybackslash}p{(\columnwidth - 4\tabcolsep) * \real{0.3333}}
  >{\raggedright\arraybackslash}p{(\columnwidth - 4\tabcolsep) * \real{0.3333}}@{}}
\toprule\noalign{}
\begin{minipage}[b]{\linewidth}\raggedright
Target
\end{minipage} & \begin{minipage}[b]{\linewidth}\raggedright
Use when
\end{minipage} & \begin{minipage}[b]{\linewidth}\raggedright
Review question
\end{minipage} \\
\midrule\noalign{}
\endhead
\bottomrule\noalign{}
\endlastfoot
\passthrough{\lstinline!STATIC!} & embed object code at link time &
which downstream libraries must still be linked? \\
\passthrough{\lstinline!SHARED!} & dynamic-link plugin or SDK
implementation & what is the visibility and ABI strategy? \\
\passthrough{\lstinline!MODULE!} & loadable runtime component & which
host loads it? \\
\passthrough{\lstinline!INTERFACE!} & header-only contract / options
aggregation & are there any build sources? \\
\passthrough{\lstinline!OBJECT!} & share objects inside one build & is
coupling to build architecture acceptable? \\
\passthrough{\lstinline!IMPORTED!} & represent externally produced
binary & is location/architecture version correct? \\
\passthrough{\lstinline!ALIAS!} & stable namespace inside current build
& does its real target own installation rules? \\
\end{longtable}

\chapter{Stage 3: Establish CTest
behavior}\label{stage-3-establish-ctest-behavior}

\begin{lstlisting}[language=sh]
ctest --preset ci-debug --output-on-failure
ctest --test-dir build/ci-debug -N
\end{lstlisting}

The reference tests assert arithmetic behavior without an external
framework. A registered test is not automatically a unit test:
classification follows the executable's dependencies and assertions.
Tests should be deterministic, independent and fast; add fixtures or
resource locks for stateful integration tests.

A failing test must return a non-zero exit code; test output alone is
not sufficient. \passthrough{\lstinline!ctest --output-junit <path>!}
can export results for CI reporting.

\chapter{Stage 4: Install and export a relocatable
package}\label{stage-4-install-and-export-a-relocatable-package}

\begin{lstlisting}[language=sh]
cmake --install build/ci-debug --prefix "$PWD/staging"
find staging -type f | sort
\end{lstlisting}

The installed layout should contain library binaries, public headers and
\passthrough{\lstinline!lib/cmake/BuildSuite!} package metadata
(platform-specific library directories may vary).
\passthrough{\lstinline!configure\_package\_config\_file!} and
\passthrough{\lstinline!write\_basic\_package\_version\_file!} generate
the package files; \passthrough{\lstinline!install(EXPORT ...)!} records
target import properties. The package should not reference the original
source directory.

\begin{lstlisting}
include(GNUInstallDirs)
include(CMakePackageConfigHelpers)
install(TARGETS bsuite_core EXPORT BuildSuiteTargets
  FILE_SET HEADERS DESTINATION "${CMAKE_INSTALL_INCLUDEDIR}"
  ARCHIVE DESTINATION "${CMAKE_INSTALL_LIBDIR}"
)
install(EXPORT BuildSuiteTargets NAMESPACE BuildSuite::
  DESTINATION "${CMAKE_INSTALL_LIBDIR}/cmake/BuildSuite")
\end{lstlisting}

\textbf{Important:} the full reference project also sets
\passthrough{\lstinline!EXPORT\_NAME core!} so
\passthrough{\lstinline!BuildSuite::core!} is the installed name,
generates the package configuration and version files, and installs
them. A consumer may only rely on the installed namespace after
\passthrough{\lstinline!find\_package!}.

\chapter{Stage 5: Build a clean external
consumer}\label{stage-5-build-a-clean-external-consumer}

\begin{lstlisting}[language=sh]
cmake -S consumer -B build/consumer -G Ninja \
  -DCMAKE_PREFIX_PATH="$PWD/staging"
cmake --build build/consumer
./build/consumer/bsuite_consumer  # adapt for Windows
\end{lstlisting}

This is a stronger contract test than linking the CLI in the producer
build: it verifies that exported targets, header installation and
package configuration are sufficient. Repeat after copying
\passthrough{\lstinline!staging!} to a different directory and
configuring a \textbf{new} consumer build against the relocated prefix.

\chapter{Stage 6: Packages, release metadata and
validation}\label{stage-6-packages-release-metadata-and-validation}

The demo's CPack configuration builds a TGZ archive. Production packages
must additionally validate manifest contents, license notices,
dependency redistribution terms, runtime loader behavior, install/remove
semantics, signing and operating-system-specific packaging requirements.

\begin{lstlisting}[language=sh]
cpack --config build/ci-release/CPackConfig.cmake -G TGZ
\end{lstlisting}

Release validation checks: clean Release configure/build/test; staging
to a fresh prefix; installed-consumer test; archive extraction and one
representative smoke test on each supported target OS.

\chapter{Stage 7: Adding third-party dependencies
safely}\label{stage-7-adding-third-party-dependencies-safely}

Prefer \passthrough{\lstinline!find\_package(SomeLib CONFIG REQUIRED)!}
for a validated installed package. Prefer imported targets such as
\passthrough{\lstinline!SomeLib::SomeLib!}; provide their minimum
versions and compatibility semantics.
\passthrough{\lstinline!FetchContent!} is useful when source integration
is intentional, but fetching/building third-party CMake files executes
their configure logic. The exact revision, checksum and provenance must
be reviewed.

\begin{lstlisting}
include(FetchContent)
FetchContent_Declare(example_dep
  GIT_REPOSITORY https://example.invalid/approved/repository.git
  GIT_TAG <pinned-immutable-commit>
)
# Example only: replace placeholder; never execute with fake location.
\end{lstlisting}

\textbf{{[}ENGINEERING EXTENSION{]} Dependency acquisition policy:}
default to pre-provisioned, verified dependencies for release builds; if
network access is required, constrain hosts, pin immutable identifiers,
record SPDX license identifiers, check advisories and verify archived
checksums. A CMake \passthrough{\lstinline!GIT\_TAG!} name alone is not
a supply-chain attestation.

\chapter{Stage 8: Diagnostics and
inspection}\label{stage-8-diagnostics-and-inspection}

\begin{lstlisting}[language=sh]
cmake --debug-find -S . -B build/ci-debug
cmake --build build/ci-debug --verbose
cmake -LAH -N build/ci-debug
ctest --test-dir build/ci-debug -VV
\end{lstlisting}

Use \passthrough{\lstinline!CMAKE\_EXPORT\_COMPILE\_COMMANDS=ON!} for
supported generators and tools such as clang-tidy. A compilation
database is not a substitute for a faithful build graph when generated
sources or complex toolchains are involved. Avoid deleting individual
cache entries at random; clean reconfigure is the default control
experiment.

\chapter{Stage 9: Portability acceptance
test}\label{stage-9-portability-acceptance-test}

\begin{longtable}[]{@{}
  >{\raggedright\arraybackslash}p{(\columnwidth - 4\tabcolsep) * \real{0.3333}}
  >{\raggedright\arraybackslash}p{(\columnwidth - 4\tabcolsep) * \real{0.3333}}
  >{\raggedright\arraybackslash}p{(\columnwidth - 4\tabcolsep) * \real{0.3333}}@{}}
\toprule\noalign{}
\begin{minipage}[b]{\linewidth}\raggedright
Case
\end{minipage} & \begin{minipage}[b]{\linewidth}\raggedright
Command strategy
\end{minipage} & \begin{minipage}[b]{\linewidth}\raggedright
Pass criterion
\end{minipage} \\
\midrule\noalign{}
\endhead
\bottomrule\noalign{}
\endlastfoot
Linux Ninja Debug & configure/build/ctest & all tests pass \\
Linux Ninja Release & separate binary dir, install & all tests pass; no
build/source paths in installed configs \\
Windows MSVC & VS multi-config,
\passthrough{\lstinline!--config Release!} & all tests pass and consumer
finds library \\
macOS Clang & Ninja or Xcode & all tests pass and staged relocation
works \\
CLI clean state & fresh checkout, no prior cache & no undeclared
downloads or local dependencies \\
\end{longtable}

\textbf{Verification boundary:} this suite's automated execution
validates the environment available during generation; it does not
demonstrate that Windows, macOS, Qt, Android or embedded targets were
tested unless their toolchains were explicitly run.

\chapter{Definition of done}\label{definition-of-done}

\begin{itemize}
\tightlist
\item[$\square$]
  Public header reachable \textbf{only through} target usage
  requirements.
\item[$\square$]
  Debug and Release use separate binary trees or an intentional
  multi-config generator.
\item[$\square$]
  CTest test fails on an induced incorrect result, then passes after
  correction.
\item[$\square$]
  Install tree contains no references to source or binary directories.
\item[$\square$]
  Independent consumer locates and links
  \passthrough{\lstinline!BuildSuite::core!}.
\item[$\square$]
  Build matrix documents all tested and untested platforms.
\item[$\square$]
  CI artifacts include machine-readable test results and immutable
  source revision.
\end{itemize}

\chapter{Further reading}\label{further-reading}

Textbook Chapters 4, 16, 19, 27, 34, 35, 40, 42. Official CMake:
\url{https://cmake.org/cmake/help/latest/guide/importing-exporting/} ;
\url{https://cmake.org/cmake/help/latest/manual/cmake-presets.7.html} .

\end{document}
